\documentclass[10pt,letterpaper]{article}
\usepackage{lmodern}
\usepackage[T1]{fontenc}
\usepackage[utf8]{inputenc}
\usepackage[letterpaper,margin=0.62in,top=0.55in,bottom=0.6in]{geometry}
\usepackage{titlesec}
\usepackage{enumitem}
\usepackage{tabularx}
\usepackage{fancyhdr}
\usepackage{lastpage}
\usepackage[hidelinks]{hyperref}

\pagestyle{fancy}
\fancyhf{}
\renewcommand{\headrulewidth}{0pt}
\fancyfoot[R]{\small Cantor \thepage/\pageref{LastPage}}
\fancypagestyle{first}{\fancyhf{}\renewcommand{\headrulewidth}{0pt}}

\setlength{\parindent}{0pt}
\hyphenpenalty=10000 \exhyphenpenalty=10000
\titleformat{\section}{\large\bfseries}{}{0pt}{}[\vspace{1pt}\titlerule]
\titlespacing*{\section}{0pt}{11pt}{5pt}
\setlist[itemize]{label=--,leftmargin=12pt,itemsep=1.5pt,topsep=2pt,parsep=0pt}

% \entry{left bold}{right}{left italic}{right italic}
% The right column is a fixed-width box, so dates and GPA can never run off the page.
\newcommand{\entry}[4]{%
  \begin{tabularx}{\textwidth}{@{}X@{\hspace{8pt}}r@{}}
    \textbf{#1} & #2 \\
    \ifx\relax#3\relax\else \textit{#3} & \textit{#4} \\ \fi
  \end{tabularx}\vspace{1pt}}

\begin{document}
\thispagestyle{first}

\begin{center}
  {\LARGE\bfseries Dylan Cantor}\\[5pt]
  Brooklyn, NY \enspace$\cdot$\enspace (610) 716-6913 \enspace$\cdot$\enspace
  \href{mailto:cantor.dylan@gmail.com}{cantor.dylan@gmail.com}\\[2pt]
  \href{https://dylancantor.com}{dylancantor.com} \enspace$\cdot$\enspace
  \href{https://www.linkedin.com/in/dylan-cantor-909b77352/}{linkedin.com/in/dylan-cantor-909b77352} \enspace$\cdot$\enspace
  \href{https://github.com/djc8582}{github.com/djc8582}
\end{center}

\section{Education}
\entry{New York University}{Expected May 2027}
  {Double major in Music Technology (Steinhardt) and Computer Science (CAS); Minor, Mathematics}{GPA 3.69}
\entry{NYU Paris --- IRCAM program}{Spring 2026}
  {Spatial audio at IRCAM, one day weekly: ambisonics, directivity, binaural rendering}{Paris, France}
\textbf{Relevant coursework.} Digital Signal Processing; Operating Systems; Artificial Intelligence; Algorithms;
Linear Algebra; Probability \& Statistics; Analog \& Digital Electronics with labs; Analog \& Digital Recording;
Electronic \& Computer Music Synthesis; MIDI Technology II (Max/MSP).

\section{Research Experience}
\entry{Research Team Member --- Georgia Tech Center for Music Technology}{May 2026 -- Present}
  {Prof.\ Gil Weinberg's research group}{Atlanta, GA (remote)}
\begin{itemize}
  \item Sole engineer on the iOS client for HearMe, a shipped application that generates a full arrangement
  from a musician's own playing, returning separable stems, chords, and time-aligned lyrics.
  \item Built a Web-Audio-style playback engine in React Native with sample-accurate scheduling for the
  piano and drum interfaces, and offline MIDI rendering of generated results.
  \item Implemented the full input path --- recording, upload, keyboard capture, lyric entry --- into the
  generation service, with a background watcher that survives app suspension.
  \item Built the result surface: one transport synchronized across stems, chords, lyrics, and play-along mode,
  with per-project version history. 200+ commits. Joined for a 150-hour summer term and have remained
  on the team since, meeting weekly with the group and separately with Prof.\ Weinberg.
\end{itemize}

\section{Professional Experience}
\entry{Cofounder and Sole Engineer --- Clubs}{February 2026 -- Present}
  {Creative-collaboration platform, live on the iOS App Store}{Brooklyn, NY}
\begin{itemize}
  \item Built and shipped a platform where projects are posted as cards, discovered by swiping, and applied to;
  accepted collaborators work in real-time channels. 1,976 users and 311 active public projects from
  213 creators across 49 schools, with 24,829 swipes, 3,596 applications and 2,262 messages in the last 30 days.
  \item Designed the AI pipelines: free-text self-descriptions parsed into interests and skills, content
  moderation, automatic project tagging, semantic search, and an email path that publishes a forwarded
  message as a card.
  \item Built a personalized discovery feed on a sawtooth-staircase ranking algorithm, batched and preloaded in
  the background so the first card paints instantly, plus real-time group and direct messaging with unread
  tracking, push notifications and universal deep links.
  \item Sole engineer across a multi-package monorepo spanning the iOS client, backend, admin console and
  marketing site; 2,300+ commits. Covered in the
  \href{https://nyunews.com}{\textit{Washington Square News}}. Cofounded with Nick Benak, who leads design and marketing.
\end{itemize}

\section{Projects}
\entry{Jamarama --- real-time multiplayer jam system; Winner, Audiotool Challenge}{June 2026}{}{}
\begin{itemize}
  \item Led the audio architecture at the Google DeepMind Magenta hackathon, with Sai Nellutla. Two to four
  players each control one dimension of a single fused track from their phones, and one player's chord
  change re-roots the whole band on the next downbeat.
  \item Placed the master clock in the host browser, fanning beats over WebSocket to phones acting as thin
  controllers; the lead keyboard fires a local tone in under 20\,ms so playing never waits on the network.
  \item Measured the generative model for tempo stability, found it drifts within roughly 30 seconds with
  no tempo input, and removed it from the clock entirely --- a deterministic clock-locked band wrapped in
  an arhythmic texture layer that floats off the grid, where drift is inaudible.
  \item Ran that layer as Magenta RealTime 2 on-device via MLX at $\sim$2.37$\times$ real-time, conditioned on
  blended player taste, live chord, and crowd energy. Tone.js, WebSocket, Node, Python, MLX.
\end{itemize}
\entry{``I Feel Like'' --- second-order ambisonic production and binaural mix}{Spring 2026}{}{}
\begin{itemize}
  \item Produced a multitrack piece entirely in the ambisonic domain at IRCAM: about forty sources
  individually encoded into a second-order field and positioned by azimuth and elevation rather than panned.
  \item Ran room encoding, directivity shaping and a feedback-delay-network reverb inside the ambisonic
  domain rather than on a stereo bus, then decoded the same scene two ways --- loudspeaker array by
  all-round ambisonic decoding, headphones by binaural decoding.
  \item Recorded all source material: field recordings on omnidirectional capsules, instruments from the
  Baschet brothers' \textit{Structures Sonores} sound sculptures, and tetrahedral-microphone material
  converted A-format to B-format before encoding. Reaper, IEM Plug-in Suite, Sennheiser AMBEO.
\end{itemize}

\section{Creative and Performance Work}
\entry{Growing --- character-based act with Miller Fishback, on Cooking Pot}{2022 -- Present}{}{}
\begin{itemize}
  \item Each alias is a character with its own register, letting the project move between rap, avant-garde
  jazz-rock and electronic music. Seven names on Spotify and Bandcamp, including \textit{rat's nest} (2024).
  \item Roughly 990,000 Spotify streams across the catalogue, including a 650,000-stream collaboration and
  released work with Open Mike Eagle, Quelle Chris and Denmark Vessey.
\end{itemize}
\entry{Close Quarters --- band with Nick Benak, on Cooking Pot}{Present}{}{}
\begin{itemize}
  \item \textit{EDUCATION}, a 49-minute continuous psychedelic rock concept album about going through school at
  NYU. Written, recorded and produced together.
\end{itemize}

\section{Technical Skills}
\textbf{Audio and DSP:} Max/MSP, RNBO, Web Audio API, JUCE, ambisonics and binaural rendering,
real-time and streaming audio, sample-accurate scheduling\\
\textbf{Machine learning:} generative audio, music information retrieval, on-device inference (MLX)\\
\textbf{Languages and frameworks:} Python, TypeScript, Swift, C++; React Native, Node, Supabase\\
\textbf{Tools:} Git, Reaper, Ableton Live, Logic Pro, IEM Plug-in Suite

\end{document}
